\section{Quasi-Distributions}

\subsection{Definition and relation to TMDs}

Since one cannot have light-like separations on the lattice,
it was proposed \cite{Ji:2013dva} to consider
spacelike separations
$z= (0,0,0,z_3)$ [or, for brevity, \mbox{$z=z_3$}]. Then,
in the \mbox{$p=(E, 0_\perp, P)$} frame,
one introduces the quasi-PDF $Q(y, P)$ through a parametrization
\begin{align}
\langle p |   \phi(0) \phi (z_3)|p \rangle
=  &
\int_{-\infty}^{\infty}   dy \,
Q(y, P) \,  e^{i y  P z_3 } \,
\  .
\label{newVDFxzQ}
\end{align}
According to this definition, the function $Q(y,p)$ characterizes the probability
that the parton carries fraction $y$ of hadron's third momentum
component $P$. Viewing the matrix element as a function
of the $\nu$ and $-z^2$ variables
(they are given by $Pz_3$ and $z_3^2$ in this case), we have
\begin{align}
{\cal M} (\nu,z_3^2)
=  &
\int_{-\infty}^{\infty}   dy \,
Q(y, P) \,  e^{i y  \nu } \,
\  .
\end{align}
Noticing that $z_3^2= \nu^2/P^2$, we get the inverse Fourier transformation
in the form
\begin{align}
Q(y,  P)   =\frac{1}{2 \pi}  \int_{-\infty}^{\infty}  d\nu \,
\, e^{-i y  \nu}  \, {\cal M} (\nu,  \nu^2/P^2)    \  .
\label{IxM}
\end{align}
It
indicates that $Q(y,  P)$ tends to $f(y)$ in the
$P \to \infty$ limit, as far as ${\cal M} (\nu,  \nu^2/P^2) \to {\cal M} (\nu, 0)$.

Thus, the deviation of the quasi-PDF $Q(y,  P)$ from the PDF $f(y)$
is determined by the dependence of ${\cal M} (\nu,  z_3^2)$
with respect to its second argument.
By \mbox{Eq. (\ref{MTMD}),} this dependence is related to the dependence
of the TMD ${\cal F} (x, \kappa^2)$ on its second argument $\kappa^2$.
Hence, the difference between $Q(y,  P)$ and $f(y)$
may be described in terms of TMDs.

To this end, we
incorporate the fact that \mbox{Eq. (\ref{MTMD}) }
is a {\it mathematical} relation between
the function \mbox{${\cal M} (\nu,  z_1^2+z_2^2) $}
and the function ${\cal F} (x, k_1^2+k_2^2)$,
no matter what is a {\it physical} meaning of the variables
$z_1,z_2$ and $k_1,k_2$.
Thus, we
substitute \mbox{Eq. (\ref{MTMD})} with $z_1=0$ and $z_2 =\nu/P$
into Eq. (\ref{IxM}) to convert it into the
expression for quasi-PDFs in terms of TMDs
\begin{align}
Q(y, P) /P =  & \,\int_{-\infty}^{\infty} d  k_1
\int_{-1} ^ {1} d  x \,   {\cal F} (x, k_1^2+(y-x)^2P^2 )
\  .
\label{QTMD}
\end{align}

Originally, this relation was derived in Ref. \cite{Radyushkin:2016hsy} using
a Nakanishi-type representation of Refs. \cite{Radyushkin:2014vla,Radyushkin:2015gpa}.
Now, we see that it is a mere consequence of Lorentz invariance.

\subsection{Quantum chromodynamics (QCD) case}

The formulas derived above are
directly applicable for non-singlet parton densities in QCD.
In that case, one deals with matrix elements of
\begin{align}
{\cal M}^\alpha  (z,p) \equiv \langle  p |  \bar \psi (0) \,
\gamma^\alpha \,  { \hat E} (0,z; A) \psi (z) | p \rangle \
\label{Malpha}
\end{align}
type, where $
{ \hat E}(0,z; A)$ is the standard $0\to z$ straight-line gauge link
in the quark (fundamental) representation.
These matrix elements may be decomposed into $p^\alpha$ and $z^\alpha$ parts:
\begin{align}
{\cal M}^\alpha  (z,p) = &2 p^\alpha  {\cal M}_p (-(zp), -z^2)
\nn & + z^\alpha  {\cal M}_z (-(zp),-z^2)
\ .
\end{align}
The ${\cal M}_p (-(zp), -z^2) $ part gives the twist-2 distribution when $z^2 \to 0$,
while $ {\cal M}_z ((zp),-z^2)$ is a purely higher-twist contamination,
and it is better to get rid of it.

If one takes $z=(z_-, z_\perp)$ in the $\alpha=+$ component of
${\cal M}^\alpha$, the $z^\alpha$-part drops out,
and one can introduce a TMD ${\cal F}(x, k_\perp^2)$ that
is related to $ {\cal M}_p (\nu, z_\perp^2)$ by the scalar formula
(\ref{MTMD}).
For quasi-distributions,
the easiest way to remove the $z^\alpha$ contamination is to take the
time component of \mbox{$ {\cal M}^\alpha  (z=z_3,p)$} and define
\begin{align}
&  {\cal M}^0   (z_3,p)
=  2p^0
\int_{-1}^1 dy\,
Q(y,P) \,
\,  e^{i  y Pz_3 }  \  .
\label{OPhixspin12}
\end{align}
Then the connection between $Q(y,P)$ and
${\cal F}(x, k_\perp^2)$ is given
by the scalar formula (\ref{QTMD}).

One may notice that the operator defining $ {\cal M}^\alpha  (z,p) $
involves a straight-line link from $0$ to $z$ rather than a
stapled link usually used in the definitions of TMDs appearing in the description of
Drell-Yan and semi-inclusive DIS processes. As is well-known, the stapled links
reflect initial or final state interactions inherent in these processes.
The ``straight-link'' TMDs, in this sense, describe the structure
of a hadron when it is in
its non-disturbed or ``primordial'' state. While it is unlikely that such a TMD can be measured
in a scattering experiment, it is a well-defined QFT object,
and one may hope that it can be measured on the lattice.

\subsection{Momentum distributions}

The quasi-PDFs describe the distribution in the fraction $y \equiv k_3/P$ of the
third component $k_3$ of the parton momentum to that of the hadron.
One can introduce distributions in $k_3$ itself:
$R (k_3,P) \equiv Q(k_3/P,P)/P$.
Then we can rewrite Eq. (\ref{QTMD}) as
\begin{align}
R(k_3, P)  =  & \,\int_{-\infty}^{\infty} d  k_1
\int_{-1} ^ {1} d  x \,   {\cal F} (x, k_1^2+(k_3-xP)^2 )
\
\label{RTMD0}
\end{align}
or, switching to the linear argument \mbox{$k_3-xP$,}
\begin{align}
R(k_3, P)  =  &  \int_{-1}^1 dx\,   {\cal R} (x, k_3-xP)
\  ,
\label{RTMD}
\end{align}
where
\begin{align}
{\cal R} (x, k_3  )   \equiv   & \,\int_{-\infty}^{\infty} d  k_1
{\cal F} (x, k_1^2+k_3^2)
\
\label{calD}
\end{align}
is the TMD $ {\cal F} (x, \kappa^2)$ integrated over the $k_1$ component
of
the two-dimensional vector $\kappa=\{k_1,k_3\}$.
According to (\ref{calD}), ${\cal R} (x, k_3  ) $ depends on $k_3$ through $k_3^2$.

For a hadron at rest, we have
\begin{align}
R(k_3, P=0)  \equiv   &\, r(k_3)=
\int_{-1}^1 dx\,   {\cal R} (x, k_3  )
\  .
\label{DTMDsmall}
\end{align}

This one-dimensional distribution may be directly obtained through a parameterization
of the density
\begin{align}
\rho (z_3^2) \equiv & {\cal M} (0,z_3^2) =
\int_{-\infty}^{\infty}   dk_3 \,
r(k_3)  \,  e^{i k_3 z_3 } \,
\
\label{dk3}
\end{align}
given by $  \langle p |   \phi(0) \phi ( z_3 )|p \rangle |_{ {\bf p} = {\bf 0} } $.
Thus, $r(k_3)$ describes a primordial distribution of $k_3$
(or any other component of ${\bf k}$) in
a rest-frame hadron.

The formula (\ref{RTMD}) has a straightforward interpretation.
According to it, when the hadron is moving, the parton's $k_3$ momentum
has two sources.

The first part, $xP$ comes from the motion
of the hadron as a whole, and the probability to get $xP$ is governed by the dependence of
the TMD ${\cal F} (x, \kappa^2)$ on its first argument, $x$.

On the other hand, the probability to get the remaining part $k_3-xP$
is governed by the dependence of the TMD on its second argument,
$\kappa^2$,
describing the primordial
rest-frame
momentum distribution.

The parameter $x$ appears in both arguments of
${\cal R} (x, k_3-xP)$ in Eq. (\ref{RTMD}),
i.e., $R(k,P)$ is given by a convolution.
In this sense, the momentum distributions $R(k,P)$
and, hence, the quasi-PDFs have a hybrid structure influenced by
the shape both of
PDFs and rest-frame distributions.

\subsection{Factorized models.}

Since the two sources of $k_3$ look like independent, it is natural
to demonstrate the hybrid nature of momentum distributions and quasi-PDFs using
a factorized model \mbox{${\cal R} (x, k_3-xP) = f(x) r (k_3-xP)$}
(the $x$ integral of $f(x)$ is normalized \mbox{to 1).}
For original ${\cal M} (\nu,-z^2)$ function, this Ansatz corresponds
to the factorization assumption ${\cal M} (\nu,-z^2) = {\cal M} (\nu,0){\cal M} (0,-z^2)$.

For illustration, we take a Gaussian form
$
\rho_G(z_3^2)=  e^{-z_3^2\Lambda^2/4}
$ for the rest-frame density. It corresponds to
\begin{align}
r_G(k_3) =  &
\frac{1}{\sqrt{\pi } \Lambda}  e^{-k_3^2/\Lambda^2}
\  .
\label{DTMDsmallfac}
\end{align}
For $f(x)$, we take
a simple PDF resembling nucleon valence densities
$f(x)=4(1-x)^3 \theta (0\leq x \leq 1)$.
As one can see from Fig. \ref{R}, the curve for $R(k,P)$
changes from a Gaussian shape for small $P$ to a shape resembling
stretched PDF for large $P$.

This result is in perfect compliance
with a known fact that wave functions
of moving hadrons are not given by a mere kinematical ``boost''
of the rest-frame wave functions.
Indeed, with increasing $P$, the impact of the
rest-frame distribution $r(k)$
is less and less visible, and eventually the shape of
$R(k,P)$ is determined by a completely
different function $f(k/P)$.

Rescaling to the $y=k/P$ variable
gives the quasi-PDF $Q(y,P)$ shown in Fig. \ref{Qgy}.
For large $P$, it clearly tends to the $f(y)$ PDF form.
In particular, using a momentum
$P \sim 10 \Lambda$ one gets a quasi-PDF that is rather close to the $P \to \infty$ limiting shape.
Still, since $\Lambda \sim \langle k_\perp \rangle$, assuming the folklore value
$\langle k_\perp \rangle \sim $ 300 MeV one translates the
\mbox{$P\sim 10 \Lambda$} estimate into $P \sim 3$ GeV,
which is uncomfortably large. Thus, a natural question is how to improve the convergence.

\begin{figure}[t]
\centerline{\includegraphics[width=3in]{images/RkP1_10_50keyn.pdf}}
\vspace{-5mm}
\caption{Momentum distributions $R(k,P)$ in the factorized Gaussian model for {$P/\Lambda =1,10,50$}.}
\label{R}
\end{figure}

\begin{figure}[t]
\centerline{\includegraphics[width=3.1in]{images/QyP1_10_50key.pdf}}
\vspace{-0.7cm}
\caption{Evolution of quasi-PDF $Q(y,P)$ in the factorized Gaussian model for {$P/\Lambda =1,10,50$}.
\label{Qgy}}
\end{figure}

\subsection{Pseudo-PDFs}

A formal reason for the complicated structure of a quasi-PDF
$Q(y,P)$
is the fact that it is obtained by the $\nu$-integral of
${\cal M} (\nu, z_3^2)e^{i \nu y}$
along a non-horizontal line $z_3=\nu/P$
in the $(\nu, z_3)$ plane (see Eq. (\ref{IxM})). With increasing $P$, its slope
decreases, the line becomes more horizontal,
and quasi-PDFs convert into PDFs.

In contrast, pseudo-PDFs ${\cal P} (x,z_3^2)$,
by definition, are given by integration of ${\cal M} (\nu, z_3^2)e^{i \nu x}$ over
horizontal lines $z_3=$ const. A very attractive feature of the pseudo-PDFs is that they
have
the $-1\leq x \leq 1$ support for all $z_3$ values. For small $z_3$,
they convert into PDFs.

More precisely, when $z_3$ is small, $1/z_3$
is analogous to the renormalization parameter $\mu$
of scale-dependent PDFs $f(x,\mu^2)$ of the standard OPE approach.

There is a subtlety, however, that while the \mbox{$\mu^2$-dependence }
of PDFs $f(x,\mu^2)$ comes solely from the evolution logarithms $\ln (\mu^2/m^2)$,
the $z_3^2$-dependence of quasi-PDFs in QCD comes both from
the evolution logarithms $\ln (z_3^2 m^2)$
and from the ultraviolet logarithms $\ln (z_3^2 \mu_R^2)$,
where $\mu_R$ is a cut-off parameter for divergences related to the gauge link
renormalization (see Ref. \cite{Craigie:1980qs}).
At the leading logarithm level, these divergences do not depend on $\nu$.
As a result, the ``reduced'' Ioffe-time distribution
\begin{align}
{\mathfrak M} (\nu, z_3^2) \equiv \frac{ {\cal M} (\nu, z_3^2)}{{\cal M} (0, z_3^2)} \
\label{redm}
\end{align}
satisfies,
for small $z_3$, the leading-order evolution equation
\begin{align}
\frac{d}{d \ln z_3^2} \,
{\mathfrak M} (\nu, z_3^2)    = - \frac{\alpha_s}{2\pi} \, C_F
\int_0^1  du \,   B ( u )  \,   {\mathfrak  M} (u \nu, z_3^2)
\label{74}
\end{align}
with respect to $1/z_3$
that
is similar to the evolution equation for $f (x,\mu^2)$ with respect to $\mu$.
The leading-order evolution kernel $B(u)$ for the non-singlet quark case
is given \cite{Braun:1994jq} by
\begin{align}
B (u)    &=  \left [ \frac{1+u^2}{1- u} \right ]_+  \
,
\label{74a}
\end{align}
with $[ \ldots ]_+$ denoting the standard ``plus'' prescription.

For the model used above (and $x \to -x$ symmetrized, as required for
non-singlet PDFs), we have
${\cal M} (\nu, 0) ={12}\left [\nu^2 - 4\sin ^2( \nu/2) \right ] /{\nu^4}    \  .
$
The shape of this function and of the convolution integral
$B \otimes {\cal M}(\nu)$ are shown in \mbox{Fig. \ref{MBM}. }
As one can see, $B \otimes {\cal M}(\nu)$ vanishes for
$\nu=0$, which reflects conservation of the vector current.
Thus, the rest-frame density ${\mathfrak M} (0,z_3^2)$ is not affected
by perturbative evolution.

\begin{figure}[t]
\centerline{\includegraphics[width=3.3in]{images/MBMkey.pdf}}
\vspace{-0.6cm}
\caption{Model Ioffe-time distribution ${\cal M} (\nu,0)$ and
the function $B \otimes {\cal M}$ governing its evolution.
\label{MBM}}
\end{figure}

\subsection{Lattice implementation}

A possible way to find the Ioffe-time distributions on the lattice
(suggested by K. Orginos) is to
calculate ${\cal M} (Pz_3, z_3^2)$ for several values of $P$,
and then to fit the results by a function of $\nu$ and $z_3^2$.

Recalling our discussion of two apparently independent
sources of obtaining $k_3$ for a moving hadron, one may hope
that
${\cal M} (\nu, z_3^2)$ factorizes, i.e.,
${\cal M} (\nu, z_3^2)= {\cal M} (\nu, 0) {\cal M} (0, z_3^2)$. Then
the reduced function ${\mathfrak M} (\nu, z_3^2) $ defined by Eq. (\ref{redm})
is equal to $ {\cal M} (\nu, 0)$,
and
the goal of obtaining ${\cal M} (\nu, 0)$ is reached.
Formally, what remains is just to take its Fourier transform to get the PDF $f(x)$.

In fact, such a factorization has been already observed
several years ago
in the pioneering study \cite{Musch:2010ka} of the
transverse momentum distributions in lattice QCD.

A serious disadvantage of quasi-PDFs is that they have the $x$-convolution
structure (\ref{QTMD}) even in a favorable situation when
the TMD [and ${\cal M} (\nu, z_3^2)$] factorizes.
On the other hand, using pseudo-PDFs in the form of
the ratio ${\mathfrak M} (\nu, z_3^2)$, one
divides out the $z_3^2$-dependence of the
primordial distribution without affecting the $\nu$-dependence
that dictates the shape of PDF.

A further advantage of using the ratio (pointed out by K. Orginos) is
the cancellation of the \mbox{$z_3$-dependence}
generated
by the lattice renormalization of
the gauge link ${ \hat E} (0,z_3; A)$.
Such a renormalization is required by
linear $|z_3| \delta m$ (where $\delta m \sim 1/a$, and $a$ is the ultraviolet cut-off)
and logarithmic $\ln (z_3^2/a^2)$ divergences \cite{Polyakov:1980ca,Dotsenko:1979wb}.
Due to their local nature, they are expected to
combine into a $\nu$-independent factor
$Z(z_3^2/a^2)$ that is the same in the numerator and denominator
of the ratio ${\mathfrak M} (\nu, z_3^2) $.

The multiplicative renormalizability of the linear
divergences of ${\cal M} (\nu, z_3^2)$
to all orders
was recently argued in Refs. \cite{Ishikawa:2016znu,Chen:2016fxx}.
A general proof for both linear and logarithmic divergences was claimed in
Ref. \cite{Ishikawa:2017faj}
on the basis of a direct analysis of relevant Feynman graphs.

Another approach \cite{Ji:2017oey,Green:2017xeu} is to treat
$ \hat E(0,z;A) $ as
$ h(0) \bar h (z) $,
where the auxiliary field $h(z)$ is analogous to the infinitely heavy
quark field of the heavy quark effective theory (HQET).
Since HQET is known to be multiplicatively renormalizable \cite{Bagan:1993zv}
this means that $\bar \psi(0)  \hat E(0,z;A) \psi (z)$ is also
multiplicatively renormalizable to all orders in perturbation theory.

In reality, ${\mathfrak M} (\nu, z_3^2)$ will have
a residual \mbox{$z_3^2$-dependence.}
It comes both
from a possible
violation of factorization for the soft part
(according to results of Ref. \cite{Musch:2010ka},
it is expected to be rather mild) and
from mandatory perturbative evolution. For a nonzero $\nu$,
the latter should be visible as a
$\ln (1/z_3^2 \Lambda^2)$ spike for small $z_3^2$.

Hence, a proposed strategy is to extrapolate
${\mathfrak M} (\nu, z_3^2)$ to $z_3^2=0$ from not too small
values of $z_3^2$, say, from those above 0.5 fm$^2$.
The resulting function ${\cal M} ^{\rm soft}(\nu,0)$
may be treated as
the Ioffe-time distribution producing the PDF $f_0(x)$
``at low normalization point''. The remaining
$\ln (1/z_3^2 \Lambda^2)$ spikes at small $z_3$ will generate its evolution.

To convert ${\mathfrak M} (\nu, z_3^2)$ into a function of $x$, one should, in principle,
know ${\mathfrak M} (\nu, z_3^2)$ for all $\nu$, which is impossible.
The maximal values of $\nu$ reached in existing lattice calculations
range from $3\pi $ \cite{Lin:2014zya} to $5\pi$ \cite{Alexandrou:2015rja} and $6 \pi $ \cite{Orginos:2017kos}.
Taking a Fourier transform in these limited ranges produces
unphysical oscillations in $x$.
Thus, the idea is to avoid the Fourier transform in $\nu$,
and just compare the reduced Ioffe-time distributions obtained from the
lattice
with those derived from experimentally known parton distributions.

Of course, an actual technical implementation of this program
should be discussed when the lattice data on ${\mathfrak M} (\nu, z_3^2)$
will become available.
